% Standalone XeLaTeX source; standard commands and distributed packages only.
% Source: public/notes/物工/問題/2021ap_exam.pdf, PDF pp. 8--9.
\documentclass[11pt,a4paper]{article}
\usepackage[margin=24mm]{geometry}
\usepackage{fontspec}
\usepackage{xeCJK}
\setCJKmainfont{HaranoAjiMincho-Regular.otf}
\setCJKsansfont{HaranoAjiGothic-Medium.otf}
\usepackage{amsmath,amssymb}
\usepackage{tikz}
\setlength{\parindent}{1em}
\setlength{\parskip}{0.35em}
\title{物理工学専攻 院試解答集 ver2\\2021年 物理学 第4問}
\author{政岡凜太郎}
\date{}
\begin{document}
\maketitle

\section*{[1.1] 電子の運動方程式}
電荷が $-q$ なので、電場から受ける力は $-q\boldsymbol E(t)$ である。復元力と減衰力を加えた Newton の式は
\[
\boxed{m\frac{d^2\boldsymbol x}{dt^2}=-q\boldsymbol E_0e^{-i\omega t}-m\omega_0^2\boldsymbol x-m\gamma\frac{d\boldsymbol x}{dt}}.
\]

\section*{[1.2] 定常応答}
$\boldsymbol x(t)=\boldsymbol x_0e^{-i\omega t}$ とすると、$\dot{\boldsymbol x}=-i\omega\boldsymbol x$、$\ddot{\boldsymbol x}=-\omega^2\boldsymbol x$ である。運動方程式に代入し、$\boldsymbol x_0$ をまとめると
\[
\left(\omega_0^2-\omega^2-i\gamma\omega\right)\boldsymbol x_0=-\frac{q}{m}\boldsymbol E_0,
\qquad
\boxed{\boldsymbol x_0=-\frac{q\boldsymbol E_0}{m\left[(\omega_0^2-\omega^2)-i\gamma\omega\right]}}.
\]
ここで減衰項の $-i\gamma\omega$ は、与えられた $e^{-i\omega t}$ 規約から決まる。

\section*{[1.3] 複素感受率}
単位体積当たりの分極振幅は
\[
\boldsymbol P_0=-Nq\boldsymbol x_0
=\frac{Nq^2}{m[(\omega_0^2-\omega^2)-i\gamma\omega]}\boldsymbol E_0
=\varepsilon_0\chi\boldsymbol E_0.
\]
$C=Nq^2/(\varepsilon_0m)$, $A=\omega_0^2-\omega^2$, $B=\gamma\omega$ とおき、分母に共役複素数を掛けると
\[
\chi=C\frac{A+iB}{A^2+B^2},\qquad
\boxed{\chi_R=\frac{Nq^2}{\varepsilon_0m}\frac{\omega_0^2-\omega^2}{(\omega_0^2-\omega^2)^2+\gamma^2\omega^2}},
\qquad
\boxed{\chi_I=\frac{Nq^2}{\varepsilon_0m}\frac{\gamma\omega}{(\omega_0^2-\omega^2)^2+\gamma^2\omega^2}}.
\]
実場の周期平均を取ると$\langle\boldsymbol E\cdot\partial_t\boldsymbol P\rangle=\varepsilon_0\omega\chi_I|\boldsymbol E_0|^2/2>0$となる。従って正の $\chi_I$ はこの時間規約で受動媒質の吸収に対応する。

\section*{[1.4] $\chi_R$ の周波数依存性}
$\omega<\omega_0$ では $\chi_R>0$、$\omega=\omega_0$ で零を通り、$\omega>\omega_0$ では負となる。共鳴近傍では正負の分散形となり、高周波では $\chi_R\to0^-$ である。図は $\gamma/\omega_0=0.12$ の例で、縦軸は正の尺度 $C/(2\gamma\omega_0)$ で規格化した概形である。
\begin{center}
\begin{tikzpicture}[x=4.5cm,y=1.15cm]
\draw[->] (0,-1.3)--(2.55,-1.3) node[right] {$\omega/\omega_0$};
\draw[->] (0,-1.35)--(0,1.35) node[above] {$\chi_R/[C/(2\gamma\omega_0)]$};
\draw[dashed] (1,-1.3)--(1,1.3);
\draw[thick,domain=0:2.4,samples=180,variable=\u]
  plot (\u,{0.24*(1-\u*\u)/((1-\u*\u)*(1-\u*\u)+0.0144*\u*\u)});
\node[below] at (1,-1.3) {$1\;(\omega=\omega_0)$};
\node[left] at (0,0) {$0$};
\end{tikzpicture}
\end{center}

\section*{[2.1] 境界での接線電場}
Faraday の法則 $\nabla\times\boldsymbol E=-\partial_t\boldsymbol B$ を、長辺が $y$ 方向で境界をまたぐ微小長方形に積分する。短辺の線積分は $\Delta z\to0$ で消え、磁束の時間変化も面積 $O(\Delta y\Delta z)$ に比例して消える。したがって上下の長辺の寄与が等しくなり、$E_y(z=0^+)=E_y(z=0^-)$ を得る。面上では共通の時間因子を省いて
\[
\boxed{E_0e^{ik\sin\theta_0 x}+E_me^{iK_{mx}x}=E_te^{iK_{tx}x}}.
\]
この関係は境界面上のすべての $x$ で成り立つ。

\section*{[2.2] 接線波数の一致}
境界に沿った位相が任意の $x$ で一致する必要があるため、入射・反射・透過波の接線成分は等しい。よって
\[
\boxed{K_{mx}=K_{tx}=k\sin\theta_0}.
\]

\section*{[2.3] Maxwell 方程式と分散関係}
自由電荷がなく一様な各媒質内では $\nabla\cdot\boldsymbol D=0$ なので $\nabla\cdot\boldsymbol E=0$ である。Faraday の法則の両辺に回転を作用させ、$\boldsymbol B=\mu_0\boldsymbol H$、$\nabla\times\boldsymbol H=\partial_t\boldsymbol D$、$\boldsymbol D=\varepsilon_0(1+\chi)\boldsymbol E$ を順に代入すると
\begin{align*}
 \nabla\times(\nabla\times\boldsymbol E)
 &=-\mu_0\frac{\partial}{\partial t}(\nabla\times\boldsymbol H)\\
 &=-\mu_0\varepsilon_0(1+\chi)\frac{\partial^2\boldsymbol E}{\partial t^2}.
\end{align*}
二重回転の公式と $\nabla\cdot\boldsymbol E=0$ より左辺は $-\nabla^2\boldsymbol E$ なので、
\[
\nabla^2\boldsymbol E=\mu_0\varepsilon_0(1+\chi)\frac{\partial^2\boldsymbol E}{\partial t^2}
\]
を得る。真空では $\chi=0$。平面波を代入すれば
\[
\boxed{k^2=\varepsilon_0\mu_0\omega^2},\qquad
\boxed{K_{mx}^2+K_{mz}^2=k^2},\qquad
\boxed{K_{tx}^2+K_{tz}^2=(1+\chi)k^2}.
\]

\section*{[2.4] 反射波の法線成分}
反射波は真空内で $+z$ 方向へ進み、その波数の大きさは $k$ である。接線成分を [2.2] に代入すると
\[
K_{mz}=+\sqrt{k^2-k^2\sin^2\theta_0}
=\boxed{k\cos\theta_0}.
\]

\section*{[2.5] 臨界角と透過波の枝}
$K_{tx}=k\sin\theta_0$ を誘電体の分散関係に代入すると
\[
K_{tz}^2=k^2(1+\chi-\sin^2\theta_0)=k^2(\chi+\cos^2\theta_0).
\]
境目 $K_{tz}=0$ は $\sin^2\theta_c=1+\chi$、すなわち
\[
\boxed{\theta_c=\arcsin\sqrt{1+\chi}=\arccos\sqrt{-\chi}}.
\]
$\theta_0<\theta_c$ では透過波は誘電体の深部 $z<0$ へ伝わる。したがって負の法線成分を選び
\[
\boxed{K_{tz}=-k\sqrt{1+\chi-\sin^2\theta_0}\quad (\theta_0<\theta_c)}.
\]
$\theta_0>\theta_c$ では $K_{tz}$ は純虚数となる。$\kappa=k\sqrt{\sin^2\theta_0-(1+\chi)}>0$ と置く。場の依存性は $e^{iK_{tz}z}$ で、媒質は $z\le0$ にあるので、$z\to-\infty$ で減衰する枝は $K_{tz}=-i\kappa$ である：
\[
\boxed{K_{tz}=-ik\sqrt{\sin^2\theta_0-(1+\chi)}\quad (\theta_0>\theta_c)}.
\]
実際 $e^{iK_{tz}z}=e^{\kappa z}=e^{-\kappa|z|}$ となる。逆の虚数符号では負の $z$ へ進むにつれて場が増大する。

\section*{[2.6] $1/e$ 減衰深さ}
振幅が境界値の $1/e$ となる深さの大きさ $z_0=|z|$ は $e^{-\kappa z_0}=e^{-1}$ より
\[
\boxed{z_0=\frac{1}{k\sqrt{\sin^2\theta_0-(1+\chi)}}
=\frac{1}{k\sqrt{-\chi-\cos^2\theta_0}}},
\qquad 0\le\cos^2\theta_0<-\chi.
\]
固定した $\chi$ では $z_0$ は $\cos^2\theta_0$ とともに増え、臨界値 $\cos^2\theta_c=-\chi$ に下側から近づくと発散する。図は無次元変数 $u=\cos^2\theta_0/(-\chi)$ を用いて形を示す。
\begin{center}
\begin{tikzpicture}[x=4.6cm,y=0.8cm]
\draw[->] (0,0)--(1.08,0) node[right] {$u=\cos^2\theta_0/(-\chi)$};
\draw[->] (0,0)--(0,4.6) node[above] {$\sqrt{-\chi}\,kz_0$};
\draw[thick,domain=0:0.94,samples=100,variable=\u]
  plot (\u,{1/sqrt(1-\u)});
\draw[dashed] (1,0)--(1,4.35);
\node[below] at (0,0) {$0$};
\node[below] at (1,0) {$1\;(\cos^2\theta_c=-\chi)$};
\node[left] at (0,1) {$1$};
\end{tikzpicture}
\end{center}
\end{document}
